\section{Exercises.sty}\label{sec:ref-exercises}

The exercise system -- the \emph{essential} package, the one the course
cannot be taught without. Small on purpose: counters, a handful of
environments, conditionals and one import wrapper, in plain LaTeX.
Self-contained: it does not depend on ProblemMeta (a stub swallows stray
\cs{ProblemMeta} blocks if that package is absent), owns its own colour and
tcolorbox loads, and works under any class. Authoring workflow:
section~\ref{sec:problem-files}; assembly workflow:
section~\ref{sec:assessments}.

\subsection*{Environments}

\begin{center}
\begin{tabularx}{\linewidth}{@{}l Y@{}}
\env{problem} & a numbered, cross-referenceable problem (flat document-wide
  counter; \cs{counterwithin} to restart per section).\\
\env{problemgroup}\texttt{\{STEM\}} & assembly-time grouping: ``Problem
  $N$.''\ $+$ stem, members render as parts (a), (b). Label the group
  \emph{inside} it, before the first part.\\
\env{qparts} / \cs{qpart} & dependent parts \emph{within} one problem file,
  lettered (a), (b); per-part points/workspace supplied at import.\\
\env{solution}\texttt{[notitle]} & printed only in a solutions build;
  the green box (\texttt{notitle} drops the corner tab). Zero-height
  overlay by default (section~\ref{sec:overlay}).\\
\env{hint} & printed only when this problem's hints are on; plain italic
  run-in under a \cs{HintLabel}. Never an overlay.\\
\env{choices}\texttt{[stacked$\vert$onepar$\vert$grid]} & lettered MC
  options; \cs{choice}\texttt{[correct]\{...\}} marks the key (boxed only
  in solutions), \cs{graphchoice}\texttt{[correct]\{FILE\}} is a graph
  option.\\
\env{matching} & two-column matching: \cs{leftitem}\texttt{[X]\{...\}} /
  \cs{leftgraph}\texttt{[X]\{FILE\}}, then \cs{rightcolumn}, then
  \cs{rightitem}\texttt{\{...\}} / \cs{rightgraph}\texttt{\{FILE\}}. The
  \texttt{[X]} names the match, shown only in solutions.\\
\env{problemchoice}\texttt{[select=, points=]} & choose-$N$-of-$M$: emits
  the instruction on every member, counts select$\times$points once.\\
\end{tabularx}
\end{center}

\subsection*{Import commands}

\cs{importproblem} wraps the \file{import} package's \cs{import} with
per-import settings (keyval family \env{eximport});
\cs{importcheckin} is the same with the check-in flavour pre-set. Bare
\cs{import} still works and yields no marker.

\begin{center}
\begin{tabularx}{\linewidth}{@{}l Y@{}}
\texttt{points=5} & the value (integer); a lone integer \texttt{[5]} is
  shorthand for it -- but \emph{only} alone: \texttt{[5, hint]} errors,
  write \texttt{[points=5, hint]}.\\
\texttt{points=2*} & announced bonus: shown, labelled, off the total.\\
\texttt{points=\{2,1*\}} & per-part list for a \env{qparts} file; star
  individual parts.\\
\texttt{pointseach=2} & the same value broadcast to every part.\\
\texttt{workspaces=\{4,4\}} & per-part handwriting space, in
  \cs{HandwritingLine}s.\\
\texttt{nocount} & shown but off the total, marker unchanged (the quiet
  variant; \texttt{*} is the announced one).\\
\texttt{hint} / \texttt{nohint} & force this file's hint on/off against
  the document toggle.\\
\texttt{number=\{2.1.8\}} & literal display number (notes): heading and
  \cs{ref} both read ``2.1.8''; no counter shown.\\
\texttt{desc=\{...\}} & parenthetical after the heading (notes).\\
\texttt{noprint} & skip the file entirely -- not read, no counter step, no
  points, no label (deliberately not typeset-then-hide).\\
\texttt{checkin} & wrap in the ``CHECK IN'' box with coloured A--E labels
  (what \cs{importcheckin} sets).\\
\end{tabularx}
\end{center}

Keys a context ignores are tolerated by design -- \texttt{points=} under
LectureNotes, \texttt{number=}/\texttt{desc=} under Assessment -- so one
import line moves between documents unchanged.

\cs{printtotalpoints} prints the two-pass sum of every \emph{counted,
printed} problem (\texttt{??} on the first compile; latexmk reruns
automatically).

\subsection*{Toggles}

The \cs{ifsolutions} and \cs{ifhints} conditionals drive the two layers,
defaulting from \cs{ShowSolutions} and \cs{ShowHints}, which build-pdfs
injects per view (section~\ref{sec:build-and-find}). Hint visibility is
three-state: the per-import keys override the document toggle.

\subsection*{Presentation hooks}

All display, no bookkeeping -- a class restyles these freely and can never
break the count (defaults in parentheses; the class overrides are in
sections \ref{sec:ref-assessment} and \ref{sec:ref-lecturenotes}):

\begin{center}
\begin{tabularx}{\linewidth}{@{}l Y@{}}
\cs{ProblemLabelWord} & the noun (\emph{Problem}).\\
\cs{ProblemNumberText} & the formatted number (\emph{Problem N}) -- one
  macro, so standalone headings and group stems always agree.\\
\cs{ProblemHeadingFormat} & a top-level problem's whole heading (run-in
  bold ``Problem N.''\ $+$ inline points).\\
\cs{ProblemGroupHeadingFormat}\texttt{\{stem\}} & the group stem's
  heading.\\
\cs{ProblemPointsFormat}\texttt{\{n\}} & how points render (inline
  ``($n$ points)'').\\
\cs{ProblemBonusPointsFormat}\texttt{\{n\}} & the bonus twin (inline
  ``($n$ points, bonus)'').\\
\cs{QPartLabelFormat} & the part-letter style (\cs{alph}; also used by
  refs, so ``N(a)'' agrees everywhere).\\
\cs{ProblemChoiceInstruction} & the choose-$N$-of-$M$ wording (``Answer
  any two of Questions 3--5.'').\\
\cs{HintLabel} & the hint run-in (\textbf{HINT:}).\\
\cs{ChoicesGridCols} & grid columns (2).\\
\cs{ChoiceLabelFormat}\texttt{\{A\}} & the option label ([A] -- square:
  round brackets are taken by parts and points, and Crowdmark bubbles are
  rectangular).\\
\cs{GraphChoiceValign} & label-to-graph alignment (\texttt{c}; \texttt{t}
  / \texttt{b}).\\
\cs{GraphLabelReserve} & width reserved beside a graph for its label
  (2.4\,em).\\
\cs{MatchColWidth}, \cs{MatchItemSep} & matching column width
  (0.38\cs{linewidth}) and item gap (2\,ex).\\
\cs{MatchKeyFormat}, \cs{CorrectChoiceFormat} & how the solutions-build
  answer cues render -- both route through the shared green
  \cs{SolutionCueBox}, so the two cues cannot drift apart
  (knobs: \cs{SolutionCueRule}, \cs{SolutionCueSep}).\\
\cs{CheckinChoiceLabel} & the check-in flashcard label (coloured A--E
  boxes; F onward falls back to plain).\\
\cs{CheckinHeaderSkip}, \cs{CheckinBodySkip} & vertical gaps inside the
  check-in box.\\
\end{tabularx}
\end{center}
